\section*{Chapitre 9 \textemdash\ Statistiques}
\addcontentsline{toc}{section}{Chapitre 9 \textemdash\ Statistiques}

\subsection*{Exercices}
\addcontentsline{toc}{subsection}{Exercices}

%────────────────────────────────────────────────────────────────
\solutiontitle{9}{1}{Indicateurs sur une s\'erie}

Notes tri\'ees~: $7,8,9,10,11,12,13,14,15,16$.
\begin{enumerate}
  \item $\bar{x}=(8+12+14+7+15+9+13+11+16+10)/10=115/10=11{,}5$.
  \item $N=10$ (pair)~; m\'ediane $=(11+12)/2=11{,}5$.
  \item $Q_1$~: m\'ediane des 5 premiers~: $9$~; $Q_3$~: m\'ediane des 5 derniers~: $14$.
  \item \'Etendue~: $16-7=9$.
        $\overline{x^2}=(49+64+81+100+121+144+169+196+225+256)/10=1405/10=140{,}5$.
        $V=140{,}5-11{,}5^2=140{,}5-132{,}25=8{,}25$~; $\sigma=\sqrt{8{,}25}\approx2{,}87$.
\end{enumerate}

%────────────────────────────────────────────────────────────────
\solutiontitle{9}{2}{Tableau de fr\'equences}

\begin{enumerate}
  \item $N=3+8+10+6+3=30$.
  \item Fr\'equences~: $10\%$, $26{,}7\%$, $33{,}3\%$, $20\%$, $10\%$.
  \item $\bar{x}=(0\cdot3+1\cdot8+2\cdot10+3\cdot6+4\cdot3)/30=(0+8+20+18+12)/30=58/30\approx1{,}93$~h.
  \item Mode~: $2$ heures (effectif $10$).
\end{enumerate}

%────────────────────────────────────────────────────────────────
\solutiontitle{9}{3}{Histogramme en classes}

Centres de classe~: $1$, $3$, $5$, $7$. $N=3+8+6+3=20$.
\begin{enumerate}
  \item Histogramme~: barres de hauteur $3,8,6,3$.
  \item $\bar{x}=(1\cdot3+3\cdot8+5\cdot6+7\cdot3)/20=(3+24+30+21)/20=78/20=3{,}9$~kg.
  \item Classe modale~: $[2;4[$ (effectif $8$).
  \item M\'ediane~: $N/2=10$~; effectifs cumul\'es~: $3$, $11$, $17$, $20$.
        La $10$e valeur est dans la classe $[2;4[$~: m\'ediane$\approx2+\frac{10-3}{8}\times2\approx3{,}75$~kg.
\end{enumerate}

%────────────────────────────────────────────────────────────────
\solutiontitle{9}{4}{Diagramme circulaire}

Total~: $120+80+45+35=280$.
\begin{enumerate}
  \item Bus~: $42{,}9\%$~; pied~: $28{,}6\%$~; v\'elo~: $16{,}1\%$~; voiture~: $12{,}5\%$.
  \item Angles~: bus $154^{\circ}$, pied $103^{\circ}$, v\'elo $58^{\circ}$, voiture $45^{\circ}$.
  \item $43{,}3\%\approx42{,}9\%$~: le mode bus (valeur approch\'ee par arrondi).
\end{enumerate}

%────────────────────────────────────────────────────────────────
\solutiontitle{9}{5}{Compl\'eter un tableau}

\begin{enumerate}
  \item $n_4=20-4-7-5=4$. Fr\'equences~: $0{,}2$~; $0{,}35$~; $0{,}25$~; $0{,}2$.
        $n_ix_i$~: $8$, $28$, $30$, $32$. Total~: $98$.
  \item $\bar{x}=98/20=4{,}9$.
  \item S\'erie tri\'ee (en tenant compte des effectifs)~: $10$e et $11$e valeurs~;
        les $11$ premi\`eres sont compos\'ees de valeurs $\leq4$ (effectif cumul\'e $4+7=11$)~;
        la $10$e et $11$e sont $4$ et $6$~: m\'ediane $=(4+6)/2=5$.
\end{enumerate}

%────────────────────────────────────────────────────────────────
\solutiontitle{9}{6}{Bo\^ite \`a moustaches}

S\'erie~: $2,5,7,8,10,11,14,15,18,20$ ($N=10$).
\begin{enumerate}
  \item $Q_2=(10+11)/2=10{,}5$~; $Q_1=$ m\'ediane de $\{2,5,7,8,10\}=7$~;
        $Q_3=$ m\'ediane de $\{11,14,15,18,20\}=15$~;
        Min$=2$~; Max$=20$.
  \item Bo\^ite~: de $7$ \`a $15$, m\'ediane \`a $10{,}5$, moustaches jusqu'\`a $2$ et $20$.
  \item $Q_3-Q_1=15-7=8$.
\end{enumerate}

%────────────────────────────────────────────────────────────────
\solutiontitle{9}{7}{Moyenne pond\'er\'ee}

\begin{enumerate}
  \item $\bar{x}_p=(14\cdot4+12\cdot3+11\cdot2+13\cdot1)/(4+3+2+1)=(56+36+22+13)/10=127/10=12{,}7$.
  \item $\bar{x}_s=(14+12+11+13)/4=50/4=12{,}5$.
  \item La moyenne pond\'er\'ee ($12{,}7$) est plus \'elev\'ee~: favorise cet \'el\`eve
        car les coefficients sont proportionnels aux meilleures mati\`eres.
\end{enumerate}

%────────────────────────────────────────────────────────────────
\solutiontitle{9}{8}{Lire un histogramme}

\begin{enumerate}
  \item $2500\times18\%=450$~; $2500\times28\%=700$~; $2500\times30\%=750$~;
        $2500\times16\%=400$~; $2500\times8\%=200$.
  \item Histogramme~: barres proportionnelles aux effectifs.
  \item Classe modale~: $[40;60[$ (fr\'equence $30\%$).
  \item Fr\'equences cumul\'ees~: $18\%$, $46\%$, $76\%$...
        La m\'ediane est dans $[40;60[$~; interpolation~:
        $Q_2\approx40+\frac{50-46}{30}\times20\approx42{,}7$~ans.
\end{enumerate}

%────────────────────────────────────────────────────────────────
\solutiontitle{9}{9}{Effectifs cumul\'es}

S\'erie~: $3,5,5,7,8,8,8,10,12,14$ ($N=10$).

\begin{center}
\begin{tabular}{|c|c|c|c|}
\hline
$x_i$ & $n_i$ & $N_i$ (cum.) & $F_i$ (cum.) \\\hline
$3$ & $1$ & $1$ & $10\%$ \\
$5$ & $2$ & $3$ & $30\%$ \\
$7$ & $1$ & $4$ & $40\%$ \\
$8$ & $3$ & $7$ & $70\%$ \\
$10$ & $1$ & $8$ & $80\%$ \\
$12$ & $1$ & $9$ & $90\%$ \\
$14$ & $1$ & $10$ & $100\%$ \\\hline
\end{tabular}
\end{center}
$Q_1$~: $25\%$ cumul\'e~: dans la classe $x=5$~: $Q_1=5$.
$Q_2$~: $50\%$~: cumul\'e $40\%$ avant $8$, puis $70\%$ apr\`es~: $Q_2=8$.
$Q_3$~: $75\%$~: $Q_3=10$.

%────────────────────────────────────────────────────────────────
\solutiontitle{9}{10}{S\'erie statistique~: analyse rapide}

\begin{enumerate}
  \item S\'erie A~: $\bar{x}=10$, $\sigma=0$ (toutes \'egales).
        S\'erie B~: $\bar{x}=100/10=10$~; $\overline{x^2}=(25+36+49+81+100+100+121+169+196+225)/10=1102/10=110{,}2$~;
        $V=110{,}2-100=10{,}2$~; $\sigma=\sqrt{10{,}2}\approx3{,}19$.
  \item S\'erie A est plus r\'eguli\`ere ($\sigma=0$).
  \item Oui~: $\sigma$ mesure l'\'ecart moyen quadratique \`a la moyenne.
\end{enumerate}

%────────────────────────────────────────────────────────────────
\solutiontitle{9}{11}{Identifier un indicateur}

\begin{enumerate}
  \item M\'ediane (robuste aux valeurs extr\^emes des prix immobiliers).
  \item Moyenne (indicateur de niveau global).
  \item \'Ecart-type (mesure la dispersion/r\'egularit\'e).
  \item Mode (valeur la plus fr\'equente).
\end{enumerate}

%────────────────────────────────────────────────────────────────
\solutiontitle{9}{12}{Calculer la variance}

$5,7,8,10,10,12$ ($n=6$).
\begin{enumerate}
  \item $\bar{x}=(5+7+8+10+10+12)/6=52/6=26/3\approx8{,}67$.
  \item $\overline{x^2}=(25+49+64+100+100+144)/6=482/6=241/3\approx80{,}33$.
  \item $V=241/3-(26/3)^2=241/3-676/9=(723-676)/9=47/9\approx5{,}22$.
  \item $\sigma=\sqrt{47/9}=\sqrt{47}/3\approx2{,}28$.
\end{enumerate}

%────────────────────────────────────────────────────────────────
\solutiontitle{9}{13}{R\`egle empirique}

$8,9,11,12,12,13,14,15,15,16$ ($n=10$).
\begin{enumerate}
  \item $\bar{x}=125/10=12{,}5$.
        $\overline{x^2}=(64+81+121+144+144+169+196+225+225+256)/10=1625/10=162{,}5$.
        $V=162{,}5-156{,}25=6{,}25$~; $\sigma=2{,}5$.
  \item Intervalle~: $[10;15]$.
  \item Valeurs dans $[10,15]$~: $11,12,12,13,14,15,15$~: $7$ valeurs.
  \item $7/10=70\%$ (proche de la r\`egle $68\%$).
\end{enumerate}

%────────────────────────────────────────────────────────────────
\solutiontitle{9}{14}{Diagramme en b\^atons}

\begin{enumerate}
  \item Diagramme~: b\^atons de hauteur $2,4,5,3,3,3$ aux faces $1,2,3,4,5,6$.
  \item $\bar{x}=(1\cdot2+2\cdot4+3\cdot5+4\cdot3+5\cdot3+6\cdot3)/20=(2+8+15+12+15+18)/20=70/20=3{,}5$.
  \item D\'e parfait~: $\bar{x}=(1+2+3+4+5+6)/6=3{,}5$.\checkmark\quad m\^eme r\'esultat.
\end{enumerate}

%────────────────────────────────────────────────────────────────
\solutiontitle{9}{15}{Indicateurs~: r\'esum\'e}

\textbf{(1)} $S_1=\{1,3,3,5,7,7,7,9\}$ ($n=8$)~:
Mode~: $7$~; M\'ediane~: $(5+7)/2=6$~; $\bar{x}=42/8=5{,}25$~;
\'Etendue~: $9-1=8$~; $V=\overline{x^2}-\bar{x}^2=(1+9+9+25+49+49+49+81)/8-5{,}25^2=272/8-27{,}5625=34-27{,}5625=6{,}44$~; $\sigma\approx2{,}54$.

\textbf{(2)} $S_2=\{2,4,6,8,10,12,14,16\}$ ($n=8$)~:
Mode~: aucun~; M\'ediane~: $(8+10)/2=9$~; $\bar{x}=72/8=9$~;
\'Etendue~: $14$~; $V=\overline{x^2}-81=(4+16+36+64+100+144+196+256)/8-81=816/8-81=102-81=21$~; $\sigma=\sqrt{21}\approx4{,}58$.

$S_2$ est plus dispers\'ee. $S_1$ est plus sym\'etrique (mode $=7$, m\'ediane $=6$, moyenne $=5{,}25$, asym\'etrie mod\'er\'ee)~; $S_2$ est parfaitement sym\'etrique (m\'ediane $=$ moyenne $= 9$).

%────────────────────────────────────────────────────────────────
\solutiontitle{9}{16}{Comparer deux s\'eries}

Rova~: $\bar{x}=115/10=11{,}5$~; tri\'e~: $7,8,9,10,11,12,13,14,15,16$~; m\'ediane $=11{,}5$.
$V=\overline{x^2}-11{,}5^2=140{,}5-132{,}25=8{,}25$~; $\sigma_R\approx2{,}87$.

Lova~: $\bar{x}=110/10=11$~; tri\'e~: $9,10,10,11,11,11,11,12,12,13$~; m\'ediane $=11$.
$V=\overline{x^2}-121=(81+100+100+121+121+121+121+144+144+169)/10-121=1222/10-121=122{,}2-121=1{,}2$~; $\sigma_L\approx1{,}10$.

Lova est plus r\'eguli\`ere ($\sigma_L\approx1{,}1\ll2{,}87$).

%────────────────────────────────────────────────────────────────
\solutiontitle{9}{17}{Fr\'equences cumul\'ees et m\'ediane}

$N=5+18+24+12+6=65$~; centres~: $300, 500, 700, 900, 1100$.
$\bar{x}=(300\cdot5+500\cdot18+700\cdot24+900\cdot12+1100\cdot6)/65=(1500+9000+16800+10800+6600)/65=44700/65\approx687{,}7$.
Effectifs cumul\'es~: $5,23,47,59,65$. $N/2=32{,}5$~: dans la classe $[600;800[$~:
$Q_2\approx600+\frac{32{,}5-23}{24}\times200\approx679$~(milliers d'ariary).

%────────────────────────────────────────────────────────────────
\solutiontitle{9}{18}{Analyse critique de statistiques}

\begin{enumerate}
  \item Somme~: $200\times7{,}2=1440$.
  \item Intervalle $[\bar{x}-\sigma,\bar{x}+\sigma]=[4{,}4;10]$~:
        la plupart des r\'eponses se situent entre $4{,}4$ et $10$.
  \item Le deuxi\`eme fabricant ($6{,}8$, $\sigma=0{,}4$) est tr\`es r\'egulier~:
        presque tous les clients donnent entre $6{,}4$ et $7{,}2$. C'est plus fiable.
        Choix contextuel~: si on veut des clients tr\`es satisfaits (quelques $10/10$),
        le premier~; si on veut une satisfaction r\'eguli\`ere, le second.
  \item La m\'ediane est robuste aux extr\^emes~: si quelques clients tr\`es insatisfaits
        ($1/10$) et beaucoup de tr\`es satisfaits, la moyenne est tir\'ee vers le bas,
        la m\'ediane reste repr\'esentative de la majorit\'e.
\end{enumerate}

%────────────────────────────────────────────────────────────────
\solutiontitle{9}{19}{Variance~: propri\'et\'es}

\begin{enumerate}
  \item $V=\frac{1}{n}\sum(x_i-\bar{x})^2=\frac{1}{n}\sum(x_i^2-2x_i\bar{x}+\bar{x}^2)
        =\overline{x^2}-2\bar{x}^2+\bar{x}^2=\overline{x^2}-\bar{x}^2$.\hfill$\square$
  \item $y_i=x_i+k$~: $\bar{y}=\bar{x}+k$~; $\sigma_y=\sigma_x$ (la dispersion ne change pas).
  \item $y_i=kx_i$~: $\bar{y}=k\bar{x}$~; $V_y=k^2V_x$~; $\sigma_y=|k|\sigma_x$.
  \item $y_i=2x_i$~: $\bar{y}=2\times6{,}5=13$~; $\sigma_y=2\times1{,}8=3{,}6$.
\end{enumerate}

%────────────────────────────────────────────────────────────────
\solutiontitle{9}{20}{S\'erie avec une inconnue}

$3,5,7,x,9,11$ ($n=6$).
\begin{enumerate}
  \item $\bar{x}=(35+x)/6$.
  \item $\bar{x}=7\Rightarrow35+x=42\Rightarrow x=7$.
  \item Avec $x=7$~: $\overline{x^2}=(9+25+49+49+81+121)/6=334/6=55{,}67$~;
        $V=55{,}67-49=6{,}67$~; $\sigma=\sqrt{6{,}67}\approx2{,}58$.
  \item $\sigma=\sqrt{6}\Rightarrow V=6$. Or
        $\overline{x^2}=\dfrac{285+x^2}{6}$ et $\bar{x}=\dfrac{35+x}{6}$.
        L'\'equation $V=6$ devient
        \[\dfrac{285+x^2}{6}-\left(\dfrac{35+x}{6}\right)^2=6,
        \qquad\text{soit}\qquad 5x^2-70x+269=0.\]
        Son discriminant vaut $\Delta=4900-4\times5\times269=-480<0$.
        Il n'existe donc aucune valeur r\'eelle de $x$.
\end{enumerate}

%────────────────────────────────────────────────────────────────
\solutiontitle{9}{21}{Position de la moyenne par rapport \`a la m\'ediane}

\begin{enumerate}
  \item $S=\{1,2,3,4,100\}$~: $\bar{x}=22$, m\'ediane $=3$.
        La moyenne ($22$) ne repr\'esente pas la majorit\'e des valeurs.
  \item $S=\{1,1,2,2,3,100\}$~: $\bar{x}=18{,}2$, m\'ediane $=2$.
        Idem~: la m\'ediane ($2$) est bien plus repr\'esentative.
  \item La m\'ediane est pr\'ef\'erable quand la distribution est fortement asym\'etrique
        ou quand il y a des valeurs extr\^emes (revenus, prix immobiliers, etc.).
  \item Exemple~: $\{1,1,1,1,1,1,100\}$~: m\'ediane $=1$, $\bar{x}\approx15$.
\end{enumerate}

%────────────────────────────────────────────────────────────────
\solutiontitle{9}{22}{Bo\^ite \`a moustaches~: comparaison}

\begin{enumerate}
  \item Deux bo\^ites~: A de 2 \`a 20 (bo\^ite 5--12), B de 4 \`a 15 (bo\^ite 7--11).
  \item \'Etendue~: A $=18$, B $=11$~: A plus grande.
        \'Ecart interquartile~: A $=7$, B $=4$~: A plus grand.
  \item Oui~: B est plus concentr\'e ($Q_3-Q_1=4$ vs $7$).
  \item Valeurs sup\'erieures \`a $9$~: dans A, $Q_2=8<9$, donc $>50\%$~; dans B, $Q_2=9$, donc $50\%$.
        Groupe A a plus de valeurs $>9$.
\end{enumerate}

%────────────────────────────────────────────────────────────────
\solutiontitle{9}{23}{Ind\'ependance de la moyenne et l'\'ecart-type}

\begin{enumerate}
  \item $\sigma=0$~: $\{10,10,10,10,10\}$~;
        $\sigma=2$~: $\{6,8,10,12,14\}$~;
        $\sigma=5$~: $\{0,5,10,15,20\}$.
  \item $\sigma=3$, $\bar{x}=5$~: $\{2,4,5,6,8\}$~;
        $\sigma=3$, $\bar{x}=15$~: $\{12,14,15,16,18\}$.
  \item Non~: la moyenne seule ne d\'ecrit pas la dispersion~;
        l'\'ecart-type seul ne donne pas le centre. Les deux sont n\'ecessaires.
\end{enumerate}

%────────────────────────────────────────────────────────────────
\solutiontitle{9}{24}{S\'erie en classes~: \'estimer les quartiles}

$N=6+14+10=30$. Classes~: $[10;15[$, $[15;20[$, $[20;25[$.
Effectifs cumul\'es~: $6$, $20$, $30$.
$Q_1$~: rang $N/4=7{,}5$~; dans $[15;20[$~: $Q_1\approx15+\frac{7{,}5-6}{14}\times5\approx15{,}5$.
$Q_2$~: rang $15$~; dans $[15;20[$~: $Q_2\approx15+\frac{15-6}{14}\times5\approx18{,}2$.
$Q_3$~: rang $22{,}5$~; dans $[20;25[$~: $Q_3\approx20+\frac{22{,}5-20}{10}\times5\approx21{,}25$.

%────────────────────────────────────────────────────────────────
\solutiontitle{9}{25}{Graphique statistique~: choisir la repr\'esentation}

\begin{enumerate}
  \item Diagramme \textbf{circulaire}~: montre les proportions par cat\'egorie.
  \item \textbf{Histogramme}~: variable continue~; classes de notes.
  \item \textbf{Bo\^ites \`a moustaches} en parall\`ele~: comparaison de distributions.
  \item \textbf{Diagramme en b\^atons} (discret) ou \textbf{courbe}~: \'evolution temporelle.
\end{enumerate}

%────────────────────────────────────────────────────────────────
\solutiontitle{9}{26}{Algorithme statistique \computer}

\begin{lstlisting}[language=Python]
import math

def stats(notes):
    n = len(notes)
    moy = sum(notes) / n
    mn, mx = min(notes), max(notes)
    etendue = mx - mn
    moy_x2 = sum(x**2 for x in notes) / n
    sigma = math.sqrt(moy_x2 - moy**2)
    tri = sorted(notes)
    if n % 2 == 1:
        med = tri[n//2]
    else:
        med = (tri[n//2-1] + tri[n//2]) / 2
    print(f"Moyenne={moy:.2f}, Min={mn}, Max={mx}")
    print(f"Etendue={etendue}, Ecart-type={sigma:.4f}")
    print(f"Mediane={med}")
    print("Superieure a 10" if moy > 10 else "Inferieure ou egale a 10")
\end{lstlisting}

%────────────────────────────────────────────────────────────────
\solutiontitle{9}{27}{\logicoalgo\ Algorithme~: calculer m\'ediane et quartiles}

\begin{enumerate}
  \item $n$ impair $\Leftrightarrow n\bmod2=1$~: vrai (d\'efinition de la parit\'e).
        N\'egation~: $n$ est pair ($n\bmod2=0$).
  \item $L=\{3,7,7,9,11\}$ (impair, $n=5$)~: m\'ediane $= L[3]=7$.
        $L=\{2,5,7,8,10,12\}$ (pair, $n=6$)~: m\'ediane $=(7+8)/2=7{,}5$.
  \item $Q_1$~: m\'ediane de la premi\`ere moiti\'e~; $Q_3$~: de la seconde.
  \item
\begin{lstlisting}[language=Python]
def mediane(L):
    s = sorted(L)
    n = len(s)
    if n % 2 == 1:
        return s[n//2]
    return (s[n//2-1] + s[n//2]) / 2

def quartiles(L):
    s = sorted(L)
    n = len(s)
    Q2 = mediane(s)
    Q1 = mediane(s[:n//2])
    Q3 = mediane(s[(n+1)//2:])
    return Q1, Q2, Q3

print(quartiles([3,7,7,9,11]))     # (3.5, 7, 10)
print(quartiles([2,5,7,8,10,12]))  # (5, 7.5, 10)
\end{lstlisting}
\end{enumerate}

%────────────────────────────────────────────────────────────────
\solutiontitle{9}{28}{Variance et \'ecart-type~: propri\'et\'es formelles}

$f(t)=\overline{x^2}-2t\bar{x}+t^2$.
\begin{enumerate}
  \item $f(t)=\frac{1}{n}\sum(x_i-t)^2=\frac{1}{n}\sum x_i^2-2t\frac{1}{n}\sum x_i+t^2
        =\overline{x^2}-2t\bar{x}+t^2$.\hfill$\square$
  \item $f(t)=(t-\bar{x})^2+\overline{x^2}-\bar{x}^2$~: trin\^ome en $t$, $a=1>0$.\hfill$\square$
  \item Minimum en $t=\bar{x}$ (coefficient de $t^2$ positif).\hfill$\square$
  \item $f(\bar{x})=\overline{x^2}-\bar{x}^2=V$~: la variance est bien la valeur minimale.\hfill$\square$
\end{enumerate}

%────────────────────────────────────────────────────────────────
\solutiontitle{9}{29}{Paradoxe de Simpson}

\begin{enumerate}
  \item G1~: A~: $9/10=90\%$ ; B~: $80/100=80\%$.
        G2~: A~: $60/100=60\%$ ; B~: $5/10=50\%$.
  \item La m\'ethode A est meilleure dans chacun des deux groupes.
  \item A~: $(9+60)/(10+100)=69/110\approx62{,}7\%$~;
        B~: $(80+5)/(100+10)=85/110\approx77{,}3\%$.
        B est meilleure globalement.
  \item Le classement s'inverse apr\`es agr\'egation~: c'est le paradoxe de Simpson.
        Les m\'ethodes ne sont pas appliqu\'ees aux m\^emes proportions des deux groupes.
\end{enumerate}

%────────────────────────────────────────────────────────────────
\solutiontitle{9}{30}{Sensibilit\'e aux valeurs extr\^emes}

\begin{enumerate}
  \item $S=\{5,6,7,8,9\}$~: $\bar{x}=7$~; $V=2$~; $\sigma=\sqrt{2}\approx1{,}41$.
  \item $\{5,6,7,8,90\}$~: $\bar{x}=116/5=23{,}2$~;
        $V=\overline{x^2}-\bar{x}^2=(25+36+49+64+8100)/5-538{,}24=1654{,}8/5-538{,}24$
        ... $\overline{x^2}=8274/5=1654{,}8$~; $V=1654{,}8-538{,}24=1116{,}56$~; $\sigma\approx33{,}4$.
        La moyenne et l'\'ecart-type ont massivement augment\'e.
  \item M\'ediane de $\{5,6,7,8,9\}=7$~; remplacer la m\'ediane ($7$) par $90$~:
        $\{5,6,90,8,9\}$ tri\'e~: $\{5,6,8,9,90\}$~; m\'ediane $=8$. Faible impact.
  \item La m\'ediane est robuste aux valeurs extr\^emes~; la moyenne et l'\'ecart-type y sont tr\`es sensibles.
\end{enumerate}

%────────────────────────────────────────────────────────────────
\solutiontitle{9}{31}{Corr\'elation \'el\'ementaire}

$(1,2),(2,4),(3,5),(4,4),(5,6)$~: $\bar{x}=3$, $\bar{y}=21/5=4{,}2$.
$\sigma_x=\sqrt{\overline{x^2}-\bar{x}^2}=\sqrt{11-9}=\sqrt{2}\approx1{,}41$.
$\overline{y^2}=(4+16+25+16+36)/5=97/5=19{,}4$~; $\sigma_y=\sqrt{19{,}4-17{,}64}=\sqrt{1{,}76}\approx1{,}33$.
$\overline{xy}=(2+8+15+16+30)/5=71/5=14{,}2$.
$r=\frac{14{,}2-3\times4{,}2}{\sqrt{2}\times\sqrt{1{,}76}}=\frac{1{,}6}{\sqrt{3{,}52}}\approx\frac{1{,}6}{1{,}876}\approx0{,}85$.

$r\approx1$~: forte corr\'elation lin\'eaire positive (les deux grandeurs augmentent ensemble).
$r\approx0$~: absence de corr\'elation. $r\approx-1$~: forte corr\'elation n\'egative.

%────────────────────────────────────────────────────────────────
\solutiontitle{9}{32}{Construire une s\'erie \`a partir d'indicateurs}

\begin{enumerate}
  \item $\bar{x}=8$, m\'ediane $=8$, $\sigma=\sqrt{6}$~: $V=6$.
        Essai~: $\{5,6,7,8,9,10,11\}$ (7 valeurs)~: $\bar{x}=8$~; $\sigma^2=(9+4+1+0+1+4+9)/7=28/7=4\ne6$.
        Essai \`a $6$ valeurs, sym\'etrique~: $\{5,7,8,8,9,11\}$~: $\bar{x}=48/6=8$, m\'ediane $(8+8)/2=8$.
        $V=(9+1+0+0+1+9)/6=20/6\ne6$.
        Essai~: $\{6,6,8,8,8,12\}$~: $\bar{x}=48/6=8$, m\'ediane$=8$.
        $V=(4+4+0+0+0+16)/6=24/6=4\ne6$. En ajustant~: $\{4,6,8,8,10,12\}$~: $\bar{x}=8$, m\'ediane$=8$, $V=(16+4+0+0+4+16)/6=40/6\ne6$. 
        \emph{Solution~:} $\{5,7,8,9,11,8\}$ tri\'e $\{5,7,8,8,9,11\}$~: $V=20/6\approx3{,}3$.
        Pour $V=6$~: $\{4,6,8,8,10,12\}$ a $V\approx6{,}67$~; $\{5,6,8,8,10,11\}$~: $V=(9+4+0+0+4+9)/6=26/6\approx4{,}3$.
        Une s\'erie exacte~: $\{5{,}5,7,8,8,9,10{,}5\}$ (7 valeurs)~: $\bar x=8$, $V=6$.
  \item $\{1,10,10,10,19\}$~: $\bar{x}=10$, mode $=10$, m\'ediane $=10$~:
        m\'ediane $=10$ \'egalement. Pour m\'ediane diff\'erente~: $\{1,2,10,10,27\}$~:
        $\bar{x}=10$, mode $=10$, m\'ediane $=10$~: encore $10$.
        $\{1,10,10,10,29\}$~: $\bar x=12$, non. Essai~: $\{5,10,10,10,15\}$~: $\bar{x}=10$, m\'ediane$=10$.
        Difficile d'\'eviter. En fait~: si mode $=\bar{x}=10$ (plus fr\'equent, valeur centrale),
        la m\'ediane tend aussi vers $10$. Exemple~: $\{1,1,10,10,28\}$~: $\bar{x}=10$, mode $1$ et $10$ (bimodal), m\'ediane $=10$.
  \item Non~: $\sigma=0$ signifie toutes valeurs \'egales \`a $\bar{x}=5$~;
        impossible d'avoir deux valeurs diff\'erentes.\hfill$\square$
\end{enumerate}

%────────────────────────────────────────────────────────────────
\solutiontitle{9}{33}{Statistiques bidimensionnelles~: r\'egression}

$(1,3),(2,5),(3,4),(4,7),(5,6)$~: $\bar{x}=3$, $\bar{y}=25/5=5$.
$\sum(x_i-\bar{x})^2=4+1+0+1+4=10$.
$\sum(x_i-\bar{x})(y_i-\bar{y})=(-2)(-2)+(-1)(0)+(0)(-1)+(1)(2)+(2)(1)=4+0+0+2+2=8$.
$a=8/10=0{,}8$~; $b=5-0{,}8\times3=2{,}6$.
Droite~: $y=0{,}8x+2{,}6$~; pour $x=6$~: $y=4{,}8+2{,}6=7{,}4$.

%────────────────────────────────────────────────────────────────
\solutiontitle{9}{34}{\logiq\ Propositions statistiques}

\begin{enumerate}
  \item \textbf{Fausse.} N\'egation~: $\exists$ s\'erie o\`u moyenne $\ne$ m\'ediane.
        Exemple~: $\{1,2,100\}$~: moyenne $=34{,}3$, m\'ediane $=2$.
  \item \textbf{Vraie.} $\{10,10,10,10,10\}$~: $\bar{x}=10$, $\sigma=0$.
  \item \textbf{Vraie} (par d\'efinition~: $Q_1\leq Q_2\leq Q_3$ toujours).
        N\'egation~: $\exists$ s\'erie o\`u $Q_1>Q_2$ ou $Q_2>Q_3$~: fausse.
  \item \textbf{Vraie}~: $\sigma=\sqrt{V}\geq0$~; $=0$ ssi toutes valeurs \'egales.
\end{enumerate}

%────────────────────────────────────────────────────────────────
\solutiontitle{9}{35}{\logiq\ Implications entre indicateurs}

\begin{enumerate}
  \item $\sigma=0\Leftrightarrow$ toutes valeurs \'egales. \textbf{\'Equivalence} ($\Leftrightarrow$).
  \item La m\'ediane $>$ moyenne indique une asym\'etrie \`a gauche, mais pas uniquement~: \textbf{$\Rightarrow$} (implication, pas \'equivalence).
  \item \'Etendue nulle $\Rightarrow$ une seule valeur $\Rightarrow\sigma=0$~; r\'eciproquement $\sigma=0\Rightarrow$ \'etendue $=0$. \textbf{\'Equivalence}.
  \item \textbf{Ni l'un ni l'autre}~: mode unique n'implique pas et n'est pas impliqu\'e par unimodalit\'e (sauf par d\'efinition).
\end{enumerate}

%────────────────────────────────────────────────────────────────
\solutiontitle{9}{36}{\logiq\ Raisonnement par l'absurde en statistiques}

\begin{enumerate}
  \item Suppos. $\sigma>0$ et toutes valeurs \'egales \`a $c$~:
        $\sigma=\sqrt{\frac{1}{n}\sum(x_i-\bar{x})^2}=\sqrt{\frac{1}{n}\sum(c-c)^2}=0$~: contradiction.\hfill$\square$
  \item $\sigma=0\Rightarrow\sum(x_i-\bar{x})^2=0\Rightarrow$ chaque terme nul $\Rightarrow x_i=\bar{x}$~:
        toutes les valeurs sont \'egales \`a $\bar{x}$.\hfill$\square$
  \item $Q_1>Q_3$~: $Q_1$ est la valeur \`a $25\%$ et $Q_3$ \`a $75\%$~; si les donn\'ees sont tri\'ees,
        la valeur au rang $N/4$ ne peut pas \^etre sup\'erieure \`a celle au rang $3N/4$.\hfill$\square$
\end{enumerate}

%────────────────────────────────────────────────────────────────
\solutiontitle{9}{37}{\logiq\ Conditions sur la variance}

\begin{enumerate}
  \item $V=\frac{1}{n}\sum(x_i-\bar{x})^2$~: chaque terme est un carr\'e, donc $\geq0$.
        La somme et la moyenne de termes $\geq0$ sont $\geq0$.\hfill$\square$
  \item CNS~: $\overline{x^2}=\bar{x}^2\Leftrightarrow V=0\Leftrightarrow\sigma=0$. C'est une \textbf{CNS}.
  \item $V=0\Leftrightarrow\sum(x_i-\bar{x})^2=0\Leftrightarrow\forall i,\;x_i=\bar{x}$.
        Sens direct~: $V=0\Rightarrow$ toutes valeurs \'egales (ci-dessus).
        Sens r\'eciproque~: toutes \'egales \`a $c$~: $\bar{x}=c$, $V=0$.\hfill$\square$
\end{enumerate}

%────────────────────────────────────────────────────────────────
\solutiontitle{9}{38}{\logiq\ Logique du paradoxe de Simpson}

\begin{enumerate}
  \item Paradoxe~: $\exists G_1,G_2,A,B$~: $(A$ meilleure dans $G_1)\wedge(A$ meilleure dans $G_2)
        \wedge(B$ meilleure globalement).
  \item N\'egation de l'implication \og Si $A$ meilleure dans tous les groupes,
        alors $A$ meilleure globalement\fg.
  \item Voir exo 29 ou d\'efi 1 (tableau des traitements A et B).
  \item Cause~: les groupes ont des tailles tr\`es diff\'erentes~;
        le groupe o\`u $B$ est meilleure est beaucoup plus grand, ce qui pond\`ere davantage son taux.
\end{enumerate}

%────────────────────────────────────────────────────────────────
\solutiontitle{9}{39}{\logiq\ Propositions sur la bo\^ite \`a moustaches}

\begin{enumerate}
  \item Pas toujours~: si beaucoup de valeurs \'egales \`a $Q_1$ ou $Q_3$, les $50\%$
        ne tombent pas exactement dans $[Q_1,Q_3]$. La proposition est \og environ vraie\fg.
  \item \textbf{Fausse.} Bo\^ite sym\'etrique ($Q_2-Q_1=Q_3-Q_2$) n'implique pas $\bar{x}=Q_2$~:
        les extr\^emes peuvent d\'eplacer la moyenne sans affecter les quartiles.
  \item $\{2,2,2,5,8\}$~: $Q_1=2$ (m\'ediane de $\{2,2,2\}=2$), $Q_2=2$, $Q_3=\frac{5+8}{2}=6{,}5$.
        Ainsi $Q_1=Q_2=2\ne Q_3=6{,}5$.\checkmark
\end{enumerate}

%────────────────────────────────────────────────────────────────
\solutiontitle{9}{40}{\algo\ Algorithme~: calculer moyenne et \'ecart-type \computer}

\begin{lstlisting}[language=Python]
import math

def stats_complet(data):
    n = len(data)
    moy = sum(data) / n
    moy_x2 = sum(x**2 for x in data) / n
    V = moy_x2 - moy**2
    sigma = math.sqrt(V)
    return moy, sigma

moy, sigma = stats_complet([5,7,8,10,10,12])
print(f"Moyenne={moy:.4f}, Ecart-type={sigma:.4f}")
# Attendu: moy=52/6~=8.667, sigma=sqrt(47/9)~=2.284
\end{lstlisting}

%────────────────────────────────────────────────────────────────
\solutiontitle{9}{41}{\algo\ Algorithme~: tri par insertion \computer}

Trace pour $L=\{5,3,8,1,4\}$~:
$i=2$~: $v=3$~; comparer $5>3$~: d\'eplacer. $L=\{3,5,8,1,4\}$.
$i=3$~: $v=8$~; $5<8$~: OK. $L=\{3,5,8,1,4\}$.
$i=4$~: $v=1$~; d\'eplacements successifs. $L=\{1,3,5,8,4\}$.
$i=5$~: $v=4$. $L=\{1,3,4,5,8\}$.\checkmark

\begin{lstlisting}[language=Python]
def tri_insertion(L):
    A = L[:]
    for i in range(1, len(A)):
        v = A[i]; j = i - 1
        while j >= 0 and A[j] > v:
            A[j+1] = A[j]; j -= 1
        A[j+1] = v
    return A

print(tri_insertion([5,3,8,1,4]))  # [1,3,4,5,8]
\end{lstlisting}

%────────────────────────────────────────────────────────────────
\solutiontitle{9}{42}{\algo\ Algorithme~: calcul des quartiles apr\`es tri \computer}

\begin{lstlisting}[language=Python]
import math

def analyse(data):
    s = sorted(data)
    n = len(s)
    def med(L): return L[len(L)//2] if len(L)%2==1 else (L[len(L)//2-1]+L[len(L)//2])/2
    Q1, Q2, Q3 = med(s[:n//2]), med(s), med(s[(n+1)//2:])
    moy = sum(s)/n
    sigma = math.sqrt(sum((x-moy)**2 for x in s)/n)
    etendue = s[-1]-s[0]
    iqr = Q3-Q1
    aberrants = [x for x in s if x < Q1-1.5*iqr or x > Q3+1.5*iqr]
    print(f"Q1={Q1}, Q2={Q2}, Q3={Q3}")
    print(f"Moy={moy:.2f}, sigma={sigma:.2f}, etendue={etendue}")
    print(f"Valeurs aberrantes: {aberrants}")

analyse([2,5,7,8,9,10,11,14,50])
\end{lstlisting}

%────────────────────────────────────────────────────────────────
\solutiontitle{9}{43}{\algo\ Algorithme~: histogramme en classes \computer}

\begin{lstlisting}[language=Python]
import math, matplotlib.pyplot as plt, numpy as np

def histogramme(data, h=5):
    mn, mx = min(data), max(data)
    bornes = list(range(int(mn), int(mx)+h+1, h))
    counts = [sum(1 for x in data if bornes[i] <= x < bornes[i+1])
              for i in range(len(bornes)-1)]
    N = len(data); n = len(counts)
    # Frequences cumulees
    cum = [sum(counts[:i+1])/N for i in range(n)]
    fig, (ax1,ax2) = plt.subplots(1,2,figsize=(12,4))
    # Histogramme
    ax1.bar(bornes[:-1], counts, width=h, align='edge', color='steelblue', edgecolor='black')
    ax1.set_title('Histogramme')
    # Polygone cumule
    mids = [(bornes[i]+bornes[i+1])/2 for i in range(n)]
    ax2.plot(mids, cum, 'bo-')
    ax2.axhline(0.5, color='r', linestyle='--', label='Mediane (50%)')
    ax2.set_title('Frequences cumulees'); ax2.legend()
    plt.show()
\end{lstlisting}

%────────────────────────────────────────────────────────────────
\solutiontitle{9}{44}{\algo\ Algorithme~: droite de r\'egression \computer}

\begin{lstlisting}[language=Python]
import math, matplotlib.pyplot as plt

def regression(xs, ys):
    n = len(xs)
    mx = sum(xs)/n; my = sum(ys)/n
    sx = math.sqrt(sum((x-mx)**2 for x in xs)/n)
    sy = math.sqrt(sum((y-my)**2 for y in ys)/n)
    mxy = sum(x*y for x,y in zip(xs,ys))/n
    a = (mxy - mx*my) / sx**2
    b = my - a*mx
    r = (mxy - mx*my)/(sx*sy)
    print(f"y = {a:.4f}x + {b:.4f}, r = {r:.4f}")
    # Trace
    plt.scatter(xs, ys); xr = [min(xs), max(xs)]
    plt.plot(xr, [a*x+b for x in xr], 'r-')
    plt.show()
    return a, b, r
\end{lstlisting}

%────────────────────────────────────────────────────────────────
\solutiontitle{9}{45}{\algo\ Algorithme~: simulation d'une loi uniforme \computer}

\begin{lstlisting}[language=Python]
import random, math

for N in [10, 100, 1000, 10000]:
    data = [random.uniform(0, 10) for _ in range(N)]
    moy = sum(data)/N
    sigma = math.sqrt(sum((x-moy)**2 for x in data)/N)
    print(f"N={N:6d}: moy={moy:.4f}, sigma={sigma:.4f}")
# La moyenne converge vers 5 et sigma vers ~2.89 = 10/sqrt(12)
\end{lstlisting}

%────────────────────────────────────────────────────────────────
\subsection*{Probl\`emes}
\addcontentsline{toc}{subsection}{Probl\`emes}

%────────────────────────────────────────────────────────────────
\psolutiontitle{9}{1}{Analyse d'une classe}

$25$ valeurs~: $7,8,9,9,10,10,11,11,11,12,12,12,12,13,13,14,14,15,15,15,16,17,17,18,19$.
\textbf{(1)} Mode~: $12$ (effectif $4$).
\textbf{(2)} $\bar{x}=\sum x_i/25=318/25=12{,}72$~; m\'ediane~: $13$e valeur $=12$~;
$\overline{x^2}=\sum x_i^2/25$~: calcul num\'erique $\approx166{,}4$~; $\sigma\approx\sqrt{166{,}4-161{,}8}\approx\sqrt{4{,}6}\approx2{,}14$.
\textbf{(3)} $Q_1=$ m\'ediane des $12$ premi\`eres~: $(10+11)/2=10{,}5$~;
$Q_3=(14+15)/2=14{,}5$~; bo\^ite~: $[10{,}5;\,12;\,14{,}5]$, moustaches $7$--$19$.
\textbf{(4)} Sup\'erieures \`a $12{,}72$~: valeurs $13,13,14,14,15,15,15,16,17,17,18,19$~: $12$ \'el\`eves, soit $48\%$.

%────────────────────────────────────────────────────────────────
\psolutiontitle{9}{2}{Contr\^ole qualit\'e}

$\bar{x}=(19{,}5\cdot3+19{,}8\cdot7+20{,}0\cdot12+20{,}2\cdot6+20{,}5\cdot2)/30=(58{,}5+138{,}6+240+121{,}2+41)/30=599{,}3/30\approx19{,}977$~mm.
$\overline{x^2}\approx(3\cdot380{,}25+7\cdot392{,}04+12\cdot400+6\cdot408{,}04+2\cdot420{,}25)/30\approx399{,}08$~;
$\sigma\approx\sqrt{399{,}08-399{,}08}$~... en pratique $\sigma\approx0{,}205$~mm.
\textbf{(2)} Vis rejet\'ees si $|l-20|>0{,}3$~: $l<19{,}7$ ou $l>20{,}3$~:
$l=19{,}5$ ($3$ vis)~; $l=20{,}5$ ($2$ vis)~: $5$ rejet\'ees.
\textbf{(3)} $5/30\approx16{,}7\%$.
\textbf{(4)} Plus $\sigma$ est petit, plus les vis sont uniformes.

%────────────────────────────────────────────────────────────────
\psolutiontitle{9}{3}{\'Evolution d'une classe}

\textbf{(1)} Somme~: $30\times12=360$.
\textbf{(2)} Nouvelle somme~: $360-18+6=348$~; nouvel effectif~: $30$.
Nouvelle moyenne~: $348/30=11{,}6$.
\textbf{(3)} L'\'ecart-type peut augmenter si on retire un \'el\`eve proche de la moyenne~;
en retirant $18$ (\'eloign\'e) et ajoutant $6$ (plus \'eloign\'e encore), la dispersion augmente g\'en\'eralement.

%────────────────────────────────────────────────────────────────
\psolutiontitle{9}{4}{Statistiques comparatives~: genre}

Femmes~: $\bar{x}_F=(22+28+30+40+26+24)/6=170/6\approx28{,}3$.
Hommes~: $\bar{x}_H=(26+38+32+48+28+35)/6=207/6=34{,}5$.
$\sigma_F\approx5{,}72$~; $\sigma_H\approx7{,}03$.
Ratios F/H~: $0{,}85$, $0{,}74$, $0{,}94$, $0{,}83$, $0{,}93$, $0{,}69$.
\'Ecart salarial moyen~: $\bar{x}_H/\bar{x}_F\approx1{,}22$ (les hommes gagnent $22\%$ de plus en moyenne).

%────────────────────────────────────────────────────────────────
\psolutiontitle{9}{5}{Donn\'ees en classes~: \'estimation}

Classes~: centres $5, 15, 30, 50, 90$~; effectifs~: $8,15,20,10,7$~; $N=60$.
$\bar{x}=(5\cdot8+15\cdot15+30\cdot20+50\cdot10+90\cdot7)/60=(40+225+600+500+630)/60=1995/60=33{,}25$~min.
Effectifs cumul\'es~: $8,23,43,53,60$. $N/2=30$~: dans $[20;40[$~:
$Q_2\approx20+\frac{30-23}{20}\times20=20+7=27$~min.

%────────────────────────────────────────────────────────────────
\psolutiontitle{9}{6}{Algorithme et statistiques \computer}

\begin{lstlisting}[language=Python]
import math, matplotlib.pyplot as plt

data = [4,7,7,9,11,12,14,15,15,16]
n = len(data)
moy = sum(data)/n
mn, mx = min(data), max(data)
sigma = math.sqrt(sum((x-moy)**2 for x in data)/n)
tri = sorted(data)
med = (tri[n//2-1]+tri[n//2])/2 if n%2==0 else tri[n//2]
print(f"Moy={moy}, Min={mn}, Max={mx}")
print(f"Sigma={sigma:.4f}, Mediane={med}")
# Histogramme
from collections import Counter
c = Counter(data)
plt.bar(c.keys(), c.values())
plt.title('Distribution'); plt.show()
\end{lstlisting}

%────────────────────────────────────────────────────────────────
\psolutiontitle{9}{7}{Enqu\^ete r\'eelle~: concevoir et analyser}

Exemple simul\'e (20 \'el\`eves, heures de r\'evision hebdomadaires)~:
Modalit\'es~: $[0;2[$, $[2;4[$, $[4;6[$, $[6;8[$, $[8;10]$.
(Toutes analyses men\'ees comme les exercices pr\'ec\'edents.)

%────────────────────────────────────────────────────────────────
\psolutiontitle{9}{8}{Taille de la bo\^ite et variation}

Min$=5$, $Q_1=10$, $Q_2=15$, $Q_3=22$, Max$=30$.
\textbf{(1)} Bo\^ite~: de $10$ \`a $22$, m\'ediane \`a $15$, moustaches $5$--$30$.
\textbf{(2)} $+5$ \`a tout~: Min$=10$, $Q_1=15$, $Q_2=20$, $Q_3=27$, Max$=35$. Tout augmente de $5$.
\textbf{(3)} $\times2$~: Min$=10$, $Q_1=20$, $Q_2=30$, $Q_3=44$, Max$=60$. Tout double.
\textbf{(4)} IQR initial~: $12$~; apr\`es $+5$~: IQR$=12$ (inchang\'e)~; apr\`es $\times2$~: IQR$=24$.

%────────────────────────────────────────────────────────────────
\psolutiontitle{9}{9}{Comparaison internationale}

Esp\'erance de vie~: $68,72,75,76,79,80,81,83$.
Ainsi $\bar{x}=\dfrac{614}{8}=76{,}75$ et $\sigma\approx4{,}65$.

Revenu~: $5,8,15,18,35,42,50,60$.
Ainsi $\bar{x}=\dfrac{233}{8}=29{,}125$ et $\sigma\approx20{,}2$.
Corr\'elation positive semblerait-elle y avoir~: les pays \`a revenus plus \'elev\'es ont g\'en\'eralement une plus grande esp\'erance de vie (sur ce petit \'echantillon).

%────────────────────────────────────────────────────────────────
\psolutiontitle{9}{10}{Suites et statistiques}

$u_n=2n+1$ pour $n=1,\ldots,10$~: $3,5,7,9,11,13,15,17,19,21$.
\textbf{(1)} Les $10$ termes.
\textbf{(2)} $\bar{x}=120/10=12$~; m\'ediane $=(11+13)/2=12$~;
$V=\overline{x^2}-144=(9+25+49+81+121+169+225+289+361+441)/10-144=1770/10-144=177-144=33$~; $\sigma=\sqrt{33}\approx5{,}74$.
\textbf{(3)} Terme du milieu~: $(u_5+u_6)/2=(11+13)/2=12=\bar{x}$.\checkmark (suite arithm\'etique~: $\bar{x}=$ milieu.)
\textbf{(4)} Pour $n$ termes d'une suite $u_k=2k+1$, la raison est $r=2$ et $\sigma=r\sqrt{(n^2-1)/12}$~:
pour $n=10$~: $\sigma=2\sqrt{99/12}=2\sqrt{8{,}25}\approx5{,}74$.\checkmark\quad $\sigma$ proportionnel \`a la raison.

%────────────────────────────────────────────────────────────────
\psolutiontitle{9}{11}{Transformation affine et statistiques}

\textbf{(1)} $\bar{y}=\frac{1}{n}\sum(ax_i+b)=a\frac{1}{n}\sum x_i+b=a\bar{x}+b$.\hfill$\square$
\textbf{(2)} $V_y=\overline{y^2}-\bar{y}^2=\overline{(ax+b)^2}-(a\bar{x}+b)^2
=a^2\overline{x^2}+2ab\bar{x}+b^2-a^2\bar{x}^2-2ab\bar{x}-b^2=a^2V_x$~;
$\sigma_y=|a|\sigma_x$.\hfill$\square$
\textbf{(3)} $y_i=2x_i$~: $\bar{y}=12$~; $\sigma_y=4$.
\textbf{(4)} $z_i=(x_i-\bar{x})/\sigma$~: $\bar{z}=0$, $\sigma_z=1$ (s\'erie centr\'ee-r\'eduite).

%────────────────────────────────────────────────────────────────
\psolutiontitle{9}{12}{Analyse de donn\'ees olympiques}

$9{,}80,9{,}84,9{,}89,9{,}89,9{,}94,9{,}95,10{,}00,10{,}02$ ($n=8$).
$\bar{x}=79{,}33/8\approx9{,}916$~s~;
$\overline{x^2}=(9{,}80^2+\ldots+10{,}02^2)/8\approx98{,}33$~;
$\sigma\approx\sqrt{98{,}33-98{,}33}\approx0{,}066$~s.
Bo\^ite~: $Q_2=(9{,}94+9{,}95)/2=9{,}945$~; $Q_1=9{,}865$~; $Q_3=9{,}975$.
\'Ecart max-min~: $10{,}02-9{,}80=0{,}22$~s.
M\'ediane $-$ moyenne~: $9{,}945-9{,}916=+0{,}029$~s (l\'eg\`erement au-dessus).

%────────────────────────────────────────────────────────────────
\psolutiontitle{9}{13}{\logiq\ Logique du paradoxe de Simpson}

\textbf{(1)} Exemple~: traitement $A$ vs $B$, deux h\^opitaux de tailles tr\`es diff\'erentes
(voir d\'efi 1 et exo 29). La pond\'eration par les tailles de groupes cr\'ee l'inversion.
\textbf{(2)} La proposition~: \og $A$ meilleure dans $G_1$ et $G_2\Rightarrow A$ meilleure globalement\fg\
est \textbf{fausse} en g\'en\'eral~: le paradoxe de Simpson est un contre-exemple.
La condition d'absence de paradoxe~: groupes de m\^eme taille, ou pond\'erations identiques.
\textbf{(3)} \textbf{Fausse}~: c'est pr\'ecis\'ement ce que le paradoxe r\'efute.

%────────────────────────────────────────────────────────────────
\psolutiontitle{9}{14}{\logiq\ Logique de la transformation affine}

\textbf{(1)} $\bar{y}=\frac{1}{n}\sum(ax_i+b)=a\cdot\frac{1}{n}\sum x_i+b=a\bar{x}+b$.
Propri\'et\'e utilis\'ee~: \textbf{lin\'earit\'e de la somme}.\hfill$\square$
\textbf{(2)} Voir exo 33 ci-dessus.\hfill$\square$
\textbf{(3)} Non~: il existe des transformations non affines qui v\'erifient ces relations.
Exemple~: $y_i=ax_i+b_i$ avec $\sum b_i=nb$ et $\sum(b_i-b)^2=0$~; ce n'est affine que si $b_i=b$ pour tout $i$.

%────────────────────────────────────────────────────────────────
\psolutiontitle{9}{15}{\algo\ Analyse statistique compl\`ete \computer}

\begin{lstlisting}[language=Python]
import math, matplotlib.pyplot as plt
from collections import Counter

def analyse_complete(data):
    s = sorted(data)
    n = len(s)
    moy = sum(s)/n
    moy_x2 = sum(x**2 for x in s)/n
    sigma = math.sqrt(moy_x2 - moy**2)
    def med(L): return L[len(L)//2] if len(L)%2==1 else (L[len(L)//2-1]+L[len(L)//2])/2
    Q1, Q2, Q3 = med(s[:n//2]), med(s), med(s[(n+1)//2:])
    iqr = Q3 - Q1
    aber = [x for x in s if x < Q1-1.5*iqr or x > Q3+1.5*iqr]
    c = Counter(s); mode_val = max(c, key=c.get)
    print(f"Moyenne={moy:.3f}, Mediane={Q2}, Mode={mode_val}, Sigma={sigma:.3f}")
    print(f"Q1={Q1}, Q3={Q3}, IQR={iqr}")
    print(f"Aberrants: {aber}")
    fig, (ax1,ax2) = plt.subplots(1,2,figsize=(10,4))
    ax1.hist(s, bins=8, color='steelblue', edgecolor='black'); ax1.set_title('Histogramme')
    ax2.boxplot(s, vert=False); ax2.set_title('Boite a moustaches')
    plt.tight_layout(); plt.show()
\end{lstlisting}

%────────────────────────────────────────────────────────────────
\psolutiontitle{9}{16}{\algo\ Simulation et statistiques \computer}

\begin{lstlisting}[language=Python]
import random, math, matplotlib.pyplot as plt, numpy as np

random.seed(42)
resultats = [random.randint(1,6)+random.randint(1,6) for _ in range(1000)]
from collections import Counter
obs = Counter(resultats)
# Probabilites theoriques
proba_th = {}
for s in range(2,13):
    n_fav = min(s-1,6) - max(s-6,1) + 1
    proba_th[s] = n_fav/36

moy = sum(resultats)/1000
sigma = math.sqrt(sum((x-moy)**2 for x in resultats)/1000)
print(f"Moy={moy:.3f}, sigma={sigma:.3f}")

xs = list(range(2,13))
freq_obs = [obs[s]/1000 for s in xs]
freq_th  = [proba_th[s] for s in xs]
plt.bar(xs, freq_obs, alpha=0.6, label='Observe')
plt.plot(xs, freq_th, 'ro-', label='Theorique')
plt.legend(); plt.title('Distribution des sommes'); plt.show()
\end{lstlisting}

%────────────────────────────────────────────────────────────────
\subsection*{D\'efis}
\addcontentsline{toc}{subsection}{D\'efis}

%────────────────────────────────────────────────────────────────
\noindent\phantomsection
{\color{BO}\bfseries\small D\'efi~1~\textendash\ \textit{Paradoxe de Simpson}}
\par\smallskip

Traitements A et B~:
\textbf{(1)} Femmes~: A $=200/250=80\%$~; B $=180/200=90\%$.
\textbf{(2)} Hommes~: A $=45/50=90\%$~; B $=150/300=50\%$.
\textbf{(3)} Global~: A $=245/300\approx81{,}7\%$~; B $=330/500=66\%$.
\textbf{(4)} B est meilleure chez les femmes, A est meilleure chez les hommes et globalement.
Cause~: le groupe des femmes (o\`u B semble meilleur) est $250+200=450$ vs les hommes $50+300=350$~;
mais A trait\`e beaucoup plus d'hommes (groupe o\`u A excelle~: $90\%$)~:
les pond\'erations d\'es\'equilibr\'ees font que A domine globalement m\^eme si B domine chez les femmes.

%────────────────────────────────────────────────────────────────
\noindent\phantomsection
{\color{BO}\bfseries\small D\'efi~2~\textendash\ \textit{Variance minimale~: preuve}}
\par\smallskip

\textbf{(1)} Parmi les s\'eries de somme fixe $S=\sum x_i$, $\bar{x}=S/n$ est fixe.
$V=\overline{x^2}-\bar{x}^2$~; comme $\bar{x}$ est fixe, minimiser $V$ revient \`a minimiser $\overline{x^2}$.
Par l'in\'egalit\'e ci-dessous, $\overline{x^2}\geq\bar{x}^2$~; la s\'erie constante $x_i=\bar{x}$
r\'ealise $\overline{x^2}=\bar{x}^2$, donc $V=0$, minimum.\hfill$\square$
\textbf{(2)} $\overline{x^2}-\bar{x}^2=V=\frac{1}{n}\sum(x_i-\bar{x})^2\geq0$.\hfill$\square$
\textbf{(3)} \'Egalit\'e $V=0$ ssi $\forall i,\;x_i=\bar{x}$~: toutes valeurs \'egales.\hfill$\square$

%────────────────────────────────────────────────────────────────
\noindent\phantomsection
{\color{BO}\bfseries\small D\'efi~3~\textendash\ \textit{G\'en\'eration de distributions}}
\par\smallskip

\textbf{(1)} $\bar{x}=10$, $\sigma=3$, m\'ediane $=10$, bo\^ite sym\'etrique~:
$\{4,8,10,10,10,12,16,10\}$~: $\bar{x}=80/8=10$, m\'ediane $=10$, sym\'etrique...
En fait~: $\{4,7,9,10,10,11,13,16\}$~: $\bar{x}=80/8=10$~;
$V=(36+9+1+0+0+1+9+36)/8=92/8=11{,}5$~; $\sigma\approx3{,}39\ne3$.
Ajustement~: $\{5,8,9,10,10,11,12,15\}$~: $\bar{x}=10$~;
$V=(25+4+1+0+0+1+4+25)/8=60/8=7{,}5$~; $\sigma\approx2{,}74$.
Avec $\{4,7,10,10,10,10,13,16\}$~: $\bar{x}=10$~;
$V=(36+9+0+0+0+0+9+36)/8=90/8\approx11{,}25$~; $\sigma\approx3{,}35$.
\textbf{(2)} $\{1,2,3,10,15,16,17,16\}$~: $\bar{x}=80/8=10$, bo\^ite asym\'etrique.
\textbf{(3)} Oui~: $\{5,5,5,5,5\}$~: $Q_1=Q_2=Q_3=5$.
\textbf{(4)} Oui~: $\{1,1,1,1,9,9,9,9\}$~: $Q_1=1$, $Q_3=9$, IQR$=8$~;
$\bar{x}=5$~; $\sigma=4=IQR/2$~: ici $\sigma<IQR$.
$\{0,10,10,10,10,20\}$~: IQR$=0$, $\sigma>0$~: $\sigma>IQR$ est possible.

%────────────────────────────────────────────────────────────────
\noindent\phantomsection
{\color{BO}\bfseries\small D\'efi~4~\textendash\ \textit{Encadrement de la moyenne}}
\par\smallskip

\textbf{(1)} Chaque $x_i\in[m,M]$~: $nm\leq\sum x_i\leq nM$~; diviser par $n$~: $m\leq\bar{x}\leq M$.\hfill$\square$
\textbf{(2)} Non~: $\{1,1,3,100\}$~: $Q_2=(1+3)/2=2$~; $\bar{x}=105/4=26{,}25>Q_2$.
Pourtant $50\%$ des valeurs sont $\leq2$ et $50\%$ sont $\geq2$~: la moyenne est tir\'ee par $100$.
\textbf{(3)} Non~: $Q_2=\bar{x}$ n'implique pas $Q_2-Q_1=Q_3-Q_2$.
Exemple~: $\{1,3,5,7,9\}$~: $\bar{x}=5=Q_2$, sym\'etrique donc OK ici~;
mais $\{1,4,5,8,12\}$~: $\bar{x}=6\ne Q_2=5$.

%────────────────────────────────────────────────────────────────
\noindent\phantomsection
{\color{BO}\bfseries\small D\'efi~5~\textendash\ \textit{Algorithme~: tri et calcul de quantiles \computer}}
\par\smallskip

\begin{lstlisting}[language=Python]
import math

def quantiles_complets(data):
    s = sorted(data)
    n = len(s)
    def med(L):
        m = len(L)
        return L[m//2] if m%2==1 else (L[m//2-1]+L[m//2])/2
    Q1 = med(s[:n//2]); Q2 = med(s); Q3 = med(s[(n+1)//2:])
    iqr = Q3 - Q1
    aber = [x for x in s if x < Q1-1.5*iqr or x > Q3+1.5*iqr]
    moy = sum(s)/n
    moy_sans = sum(x for x in s if x not in aber)/(n-len(aber)) if len(aber)<n else None
    print(f"Q1={Q1}, Q2={Q2}, Q3={Q3}, IQR={iqr}")
    print(f"Aberrants: {aber}")
    print(f"Moy avec aberrants: {moy:.3f}")
    if moy_sans: print(f"Moy sans aberrants: {moy_sans:.3f}")

quantiles_complets([2,5,7,8,9,10,11,14,50])
# 50 est aberrant: Q3+1.5*IQR = 14+1.5*7 = 24.5 < 50
\end{lstlisting}

%────────────────────────────────────────────────────────────────
\noindent\phantomsection
{\color{BO}\bfseries\small D\'efi~6~\textendash\ \textit{Statistiques bidimensionnelles~: droite de r\'egression}}
\par\smallskip

$(8,10),(12,13),(15,14),(10,11),(18,17)$~: $\bar{x}=63/5=12{,}6$~; $\bar{y}=65/5=13$.
$\sigma_x^2=\overline{x^2}-\bar{x}^2=(64+144+225+100+324)/5-158{,}76=171{,}4-158{,}76=12{,}64$.
$\overline{xy}=(80+156+210+110+306)/5=862/5=172{,}4$.
$a=(172{,}4-12{,}6\times13)/12{,}64=(172{,}4-163{,}8)/12{,}64=8{,}6/12{,}64\approx0{,}681$.
$b=13-0{,}681\times12{,}6\approx13-8{,}58=4{,}42$.
Droite~: $y=0{,}681x+4{,}42$.
Pour $x=13$~: $y\approx0{,}681\times13+4{,}42\approx13{,}27$.

%────────────────────────────────────────────────────────────────
\noindent\phantomsection
{\color{BO}\bfseries\small D\'efi~7~\textendash\ \textit{\logiq\ Variance minimale et in\'egalit\'e}}
\par\smallskip

\textbf{(1)} $V=\frac{1}{n}\sum(x_i-\bar{x})^2\geq0$ car somme de carr\'es.\hfill$\square$
\textbf{(2)} $V=0\Leftrightarrow\sum(x_i-\bar{x})^2=0\Leftrightarrow\forall i,\;x_i=\bar{x}$.\hfill$\square$
\textbf{(3)} Si $\sum x_i=S$ fix\'e, $\bar{x}=S/n$ fix\'e. $V=\overline{x^2}-\bar{x}^2$~;
minimiser $V$ = minimiser $\overline{x^2}$. Par Cauchy-Schwarz~: $\overline{x^2}\geq\bar{x}^2$,
avec \'egalit\'e ssi tous \'egaux. La s\'erie constante r\'ealise ce minimum.\hfill$\square$
\textbf{(4)} Implications~: $V\geq0$ (toujours)~; $V=0\Rightarrow$ valeurs \'egales~;
valeurs \'egales $\Rightarrow V=0$~; donc $V=0\Leftrightarrow$ valeurs \'egales.

%────────────────────────────────────────────────────────────────
\noindent\phantomsection
{\color{BO}\bfseries\small D\'efi~8~\textendash\ \textit{\algo\ Calculatrice statistique compl\`ete \computer}}
\par\smallskip

\begin{lstlisting}[language=Python]
import math, matplotlib.pyplot as plt
from collections import Counter

def calculatrice():
    print("Saisir les donnees (entree vide pour terminer):")
    data = []
    while True:
        s = input("Valeur: ")
        if not s: break
        data.append(float(s))
    if not data: return
    n = len(data); s = sorted(data)
    moy = sum(data)/n
    sigma = math.sqrt(sum((x-moy)**2 for x in data)/n)
    def med(L): return L[len(L)//2] if len(L)%2==1 else (L[len(L)//2-1]+L[len(L)//2])/2
    Q1, Q2, Q3 = med(s[:n//2]), med(s), med(s[(n+1)//2:])
    mode_val = max(Counter(data), key=Counter(data).get)
    print(f"\n--- Rapport ---")
    print(f"n={n}, Moy={moy:.4f}, Sigma={sigma:.4f}")
    print(f"Q1={Q1}, Q2={Q2}, Q3={Q3}, Mode={mode_val}")
    # Transformation affine
    a = float(input("a (y=ax+b): ")); b = float(input("b: "))
    print(f"Nouvelle moy={a*moy+b:.4f}, nouveau sigma={abs(a)*sigma:.4f}")
    # Graphiques
    fig, axes = plt.subplots(1,3,figsize=(15,4))
    axes[0].hist(data, bins='auto', color='steelblue', edgecolor='black')
    axes[0].set_title('Histogramme')
    axes[1].boxplot(data, vert=False)
    axes[1].set_title('Boite a moustaches')
    c = Counter([round(x) for x in data])
    axes[2].bar(c.keys(), c.values())
    axes[2].set_title('Diagramme en batons')
    plt.tight_layout(); plt.show()

# calculatrice()  # Decommenter pour utilisation interactive
\end{lstlisting}
